%% chapters/03-dasar-teori.tex
%% Bab III — Dasar Teori
%% Aturan: lihat .cursor/rules/08-bab3-dasar-teori.mdc
%% Berisi formulasi matematis lengkap; Bab IV merujuk persamaan-persamaan di sini.

\chapter{Dasar Teori}
\label{bab:dasar-teori}

Bab ini memuat fondasi teoretis yang menjadi dasar metodologi pada
Bab~\ref{bab:metodologi}. Disajikan formulasi matematis dari prediksi
RUL, varian arsitektur LSTM (sLSTM, mLSTM), \emph{state-space model}
dan Mamba, \emph{attention}, \emph{Sparse Autoencoder}, metode
atribusi, serta metrik evaluasi.

\section{Formulasi Prediksi RUL}
\label{subsec:bab3_rul}

\subsection{Definisi RUL}
\label{subsec:bab3_def_rul}

Diberikan deret waktu sinyal getaran $\{\mathbf{x}_t\}_{t=1}^{T}$
dengan $\mathbf{x}_t \in \mathbb{R}^{d_{\text{in}}}$, RUL pada waktu
$t$ didefinisikan sebagai
\begin{equation}
y_t = T_{\text{EOL}} - t,
\label{eq:bab3_rul_def}
\end{equation}
dengan $T_{\text{EOL}}$ adalah waktu \emph{end-of-life} bantalan.
Dalam praktik, RUL dinormalisasi ke selang $[0,1]$.

\subsection{Skema Pelabelan}
\label{subsec:bab3_label}

% Linear, piecewise (constant healthy phase + linear degrading phase),
% exponential. Disertasi ini akan menggunakan pendekatan piecewise
% (lihat Bab IV).
\dots

\section{Long Short-Term Memory dan Variannya}
\label{subsec:bab3_lstm}

\subsection{LSTM Klasik}
\label{subsec:bab3_lstm_classic}

% Persamaan gating LSTM original (Hochreiter & Schmidhuber, 1997).
\begin{align}
\mathbf{f}_t &= \sigma(\mathbf{W}_f \mathbf{x}_t + \mathbf{U}_f \mathbf{h}_{t-1} + \mathbf{b}_f) \label{eq:bab3_lstm_f} \\
\mathbf{i}_t &= \sigma(\mathbf{W}_i \mathbf{x}_t + \mathbf{U}_i \mathbf{h}_{t-1} + \mathbf{b}_i) \\
\mathbf{o}_t &= \sigma(\mathbf{W}_o \mathbf{x}_t + \mathbf{U}_o \mathbf{h}_{t-1} + \mathbf{b}_o) \\
\tilde{\mathbf{c}}_t &= \tanh(\mathbf{W}_c \mathbf{x}_t + \mathbf{U}_c \mathbf{h}_{t-1} + \mathbf{b}_c) \\
\mathbf{c}_t &= \mathbf{f}_t \odot \mathbf{c}_{t-1} + \mathbf{i}_t \odot \tilde{\mathbf{c}}_t \\
\mathbf{h}_t &= \mathbf{o}_t \odot \tanh(\mathbf{c}_t)
\end{align}

\subsection{sLSTM (\emph{scalar} LSTM)}
\label{subsec:bab3_slstm}

% Pekerjaan Beck dkk. (2024). Stabilizer + exponential gates.
\dots

\subsection{mLSTM (\emph{matrix} LSTM)}
\label{subsec:bab3_mlstm}

% Persamaan mLSTM dengan matrix memory.
\begin{equation}
\mathbf{C}_t = \mathbf{f}_t \mathbf{C}_{t-1} + \mathbf{i}_t \mathbf{v}_t \mathbf{k}_t^{\top}
\label{eq:bab3_mlstm_C}
\end{equation}
\dots

\section{\emph{State-Space Models} dan Mamba}
\label{subsec:bab3_ssm}

\subsection{SSM Kontinu}
\label{subsec:bab3_ssm_kontinu}

\begin{align}
\dot{\mathbf{h}}(t) &= \mathbf{A} \mathbf{h}(t) + \mathbf{B} x(t) \\
y(t) &= \mathbf{C} \mathbf{h}(t)
\label{eq:bab3_ssm_kontinu}
\end{align}

\subsection{Diskretisasi Zero-Order Hold}
\label{subsec:bab3_ssm_diskret}

\begin{align}
\bar{\mathbf{A}} &= \exp(\Delta \mathbf{A}) \\
\bar{\mathbf{B}} &= (\Delta \mathbf{A})^{-1}(\exp(\Delta \mathbf{A}) - \mathbf{I}) \cdot \Delta \mathbf{B}
\label{eq:bab3_zoh}
\end{align}

\subsection{Mamba: Selektivitas dan \emph{Selective Scan}}
\label{subsec:bab3_mamba}

% Inovasi Mamba: $\mathbf{B}, \mathbf{C}, \Delta$ tergantung input,
% sehingga selektif. Selective scan algorithm.
\dots

\section{Mekanisme \emph{Attention}}
\label{subsec:bab3_attn}

\begin{equation}
\text{Attn}(\mathbf{Q},\mathbf{K},\mathbf{V}) = \mathrm{softmax}\!\left(\frac{\mathbf{Q}\mathbf{K}^{\top}}{\sqrt{d_k}}\right)\mathbf{V}
\label{eq:bab3_attn}
\end{equation}

\section{\emph{Sparse Autoencoder}}
\label{subsec:bab3_sae}

\subsection{Formulasi Dasar}
\label{subsec:bab3_sae_basic}

\begin{align}
\mathbf{z}_t &= \text{ReLU}(\mathbf{W}_e \mathbf{h}_t + \mathbf{b}_e) \\
\hat{\mathbf{h}}_t &= \mathbf{W}_d \mathbf{z}_t + \mathbf{b}_d \\
\mathcal{L}_{\text{SAE}} &= \|\mathbf{h}_t - \hat{\mathbf{h}}_t\|_2^2 + \lambda \cdot \Omega(\mathbf{z}_t)
\label{eq:bab3_sae_loss}
\end{align}

\subsection{Top-$k$ Sparsity}
\label{subsec:bab3_topk}

\begin{equation}
\mathbf{z}_t^{(k)} = \mathrm{TopK}_k(\mathbf{z}_t)
\label{eq:bab3_topk}
\end{equation}

\section{Metode Atribusi: SHAP dan \emph{Integrated Gradients}}
\label{subsec:bab3_attribution}

\subsection{SHAP}
\label{subsec:bab3_shap}
\dots

\subsection{Integrated Gradients}
\label{subsec:bab3_ig}

\begin{equation}
\mathrm{IG}_i(\mathbf{x}) = (x_i - x'_i) \int_0^1 \frac{\partial f(\mathbf{x}' + \alpha(\mathbf{x}-\mathbf{x}'))}{\partial x_i}\, d\alpha
\label{eq:bab3_ig}
\end{equation}

\section{Frekuensi Karakteristik Bantalan}
\label{subsec:bab3_freq_char}

% Persamaan BPFO, BPFI, BSF, FTF. Lihat .cursor/rules/15-domain-rul-bearings.mdc.
\dots

\section{Metrik Evaluasi}
\label{subsec:bab3_metrik}

\begin{align}
\text{RMSE} &= \sqrt{\frac{1}{N}\sum_{i=1}^{N}(y_i - \hat{y}_i)^2} \label{eq:bab3_rmse} \\
\text{MAE}  &= \frac{1}{N}\sum_{i=1}^{N}|y_i - \hat{y}_i| \label{eq:bab3_mae}
\end{align}

\section{Notasi yang Digunakan}
\label{subsec:bab3_notasi}

% Tabel notasi terpadu — wajib ada untuk konsistensi lintas bab.
\begin{table}[ht]
\centering
\caption{Notasi terpadu yang digunakan dalam disertasi.}
\label{tab:bab3_notasi}
\small
\begin{tabular}{@{}p{2.5cm}p{11cm}@{}}
\toprule
Notasi & Arti \\
\midrule
$\mathbf{x}_t$ & Vektor input pada waktu $t$, $\mathbf{x}_t \in \mathbb{R}^{d_{\text{in}}}$ \\
$\mathbf{h}_t$ & Vektor laten / hidden state, $\mathbf{h}_t \in \mathbb{R}^{d_{\text{model}}}$ \\
$\mathbf{C}_t$ & Matriks memori mLSTM, $\mathbf{C}_t \in \mathbb{R}^{d \times d}$ \\
$\Delta$       & Step size diskretisasi SSM \\
$\mathbf{z}_t$ & Aktivasi laten SAE \\
$\lambda, k$   & Parameter regularisasi dan Top-$k$ SAE \\
\bottomrule
\end{tabular}
\end{table}
